\section{本章小结}
\label{sec:graph-objects-summary}

本章开头问：图记录了什么，沿连接能到哪里，删除连接会怎样，结构又怎样记录和比较。第一个问题的回答是，图$\mathcal G=(V,E)$只记录有限个对象之间的直接连接，与纸上的画法无关，比较两幅画只需比较顶点集与边集。方向、边权和属性各自附加一类信息。方向决定哪些顶点序列可以通行，忽略方向得到的底层无向图丢掉了先后。边权的含义由问题决定：长度沿路线相加，容量限制单条边的通过量，相似度鼓励两端接近；零权边仍是边，“能否到达”查边集，“累计多少”查边权。子图可以任意保留一部分边，诱导子图则保留选定顶点之间的全部边，因而由顶点集唯一确定。

沿连接能到哪里，由行走和路径回答。行走允许重复，便于拼接与计数；路径不重复顶点，对应实际路线。任何行走都可以随读随删地化成一条同端点的路径，所以用两者定义的可达关系相同。无向图中可达关系是等价关系，它把顶点集唯一地划分为连通分量；核对一个声称的划分，只需检查没有边跨越不同的块、每块内部连通。有向图中可达不对称，要区分要求双向到达的强连通与忽略方向的弱连通，前者蕴含后者，反之不然。任务能否排出先后，取决于有没有顺着箭头绕回起点的有向环；不含有向环的有向无环图，底层无向图仍可以有环。

删除连接的影响由分隔集、割点和桥刻画，判断的依据始终是有没有替代路径。核对分隔集，就是在删除后的图中求出可达集；一条边是桥当且仅当它不在任何环上，删边至多使分量数增加一，删点却可能把分量拆成许多块。用最少的边保持连通，就得到树：树等价于任意两点之间路径唯一，$n$个顶点的树恰有$n-1$条边。连通、无环、恰有$n-1$条边三个条件中任意两个推出第三个，所以判断树只需数边再查连通。选根把树组织成父子层次而不改变边集；生成森林在每个分量上保留一棵生成树，恰有$n-c$条边，并且一般不唯一，它丢掉了原图的替代路径。

同一张图可以用边列表、邻接表、邻接矩阵和关联矩阵记录。四者保存的连接相同，支持的访问和代价不同，转换本身也有代价，选择时应先列出算法反复执行的操作。邻接矩阵还把组合计数变成矩阵运算：$\bm A^k$的条目计数长度为$k$的行走，其中包含重复顶点，不能直接当作路径数。

同构用保持邻接与不相邻的双射刻画“只差一套名称”，在矩阵上表现为置换共轭$\bm A'=\bm P\bm A\bm P^{\top}$。给定对应时，核对边的像并比较边数就够了。度、分量数、特征值等不变量可以排除同构，却不足以确认同构，而逐个尝试$n!$个对应并不可行。

本章说明了怎样在小图上核对连通分量、强连通分量、生成树、无环性和同构，却还没有给出在大图上高效求出它们的方法。第~\ref{chap:graph-traversal}章在邻接表上用广度优先和深度优先遍历实际计算这些对象，用拓扑排序安排有向无环图的次序，并用顶点标记细化比较图的结构；那里的算法输出，仍要按本章的定义与核对方法检验。
